\section{Configuración del Entorno de Desarrollo}
La Configuración del entorno de desarrollo consisten principalmente en la instalación y configuración
del software a utilizar, siendo estos Blender y Python.

\subsubsection{Configuración del IDE}
Se realizó una instalación estandar del editor de Visual Studio Code, agregando las extensiones de 
Python, Pylance, Python Debugger y Python Environment para compatibilidad con Python y sus herramientas
de desarrollo, así también como las extensiones de Remote Repositories y GitHub pull requests para la
administración de repositorios remotos de GitHub dentro del editor. 

\subsubsection{Configuración de Python}
Para el desarrollo del proyecto se utilizarán dos versiones de Python. La versión 3.13.13 se instala
por ser la misma versión utilizada por Blender 5.2 LTS, mientras que se instala la versión 3.12
por ofrecer el mejor equilibrio entre compatibilidad y modernidad para utilizar Pytorch y otras
bibliotecas utilizadas para el componente de extracción. Para ambas versiones se realiza una instalación
estandar en el sistema y se indica la instalación específica para cada ambiente.

\subsubsection{Ambientes virtuales}
Se crea un ambiente virtual para cada componente desarrollado. Ambos ambientes son creados bajo el diretorio \textit{.venv} 
como es convención, omitiendo la instalación de paqueterías por defecto. 
Los ambientes son administrados a través de las extensiones de Python para Visual Studio Code.
\medskip

El ambiente virtual del componente
de extracción se crea utilizando Python 3.12 como el binario de ejecución, bajo el nombre \textit{venv-extraccion}.
Bajo este ambiente se instalaron las librerias torch (pytorch), numpy, y openvdb.
\medskip

El ambiente virtual del componente de visualización se crea con Python 3.13, bajo el nombre \textit{venv-visualización}.
Bajo este ambiente se instaló la libreria de bpy y openvdb. Este ambiente es meramente de desarrollo,
la ejecución de los scripts generados será a través del interprete del ambiente integrado de
Python en Blender.
\medskip

Cabe hacer mención que la distribución oficial de la biblioteca de OpenVDB no ha sido actualizada desde la
versión 3.8 de Python, por lo que se optó por compilar manualmente para las versiones 3.12 y 3.13. Las
instrucciones de compilación e instalación en el ambiente, junto con los archivos generados y bibliotecas
requeridas, son incluidos como anexos.

\subsubsection{Blender}
Se realizó una instalación estandar de Blender 5.2. En las preferencias de Blender se configuró para que
utilizar CUDA como backend de renderizado de Cycles y la tarjeta gráfica dedicada del sistema como dispositivo 
de renderizado. Se configuró Vulkan como backend de gráficos, utilizando también la tarjeta gráfica dedicada.
\medskip

Adicionalmente se configuró la ubicación del código fuente del componente de visualización como directorio adicional 
de scripts para facilitar su descubrimiento y ejecución por parte de Blender.